\chapter{Theoretical Background}\label{ch:theory}
This chapter establishes the theoretical foundation for the research by connecting the macroeconomic concepts of \gls{mp} with advanced language technology. First, it explores the evolution of central bank communication, emphasising the structural nuances of the \gls{ecb}'s press conferences. Second, it introduces Natural Language Processing as a methodological bridge, tracing the shift from traditional sentiment lexicons to specialised transformer models capable of decoding \gls{ecb} narratives. Finally, the chapter defines the quantitative targets ranging from official policy rates to high-frequency market expectations and structural shocks, providing the framework needed to assess the real-world financial impact of central bank text.

\section{Central Bank Communication}
Over the past three decades, central banking has undergone fundamental changes in its disclosure structures. Before the 1990s, there was no public explanation of why monetary policy decisions were made. Policymakers believed that a degree of unpredictability enhanced the effectiveness of policy interventions. This period of institutional unpredictability was characterised by a marked lack of formal public guidance \cite[p. 3]{NBERw11898}. Market participants repeatedly had to guess the target interest rate and the bank's policy setting. They relied on the size, timing, and nature of the operations conducted by the \gls{fed} trading departments in the open market.

Since then, communication has evolved from a peripheral public relations function to a core instrument of \gls{mp}. Central banks have increasingly recognised transparency and accountability as fundamental public service obligations \cite[p. 15--18]{blinder2008central}. Beyond institutional accountability, the evolution of \gls{mp} frameworks has demonstrated that transparent and credible central banks are considered important for maintaining long-term price stability \cite{daly2025dynamic}. This is particularly evident in the widespread adoption of inflation targeting across advanced and emerging economies. Empirically, clear communication has been shown to enhance central bank independence, strengthen institutional credibility, and anchor inflation expectations in complex macroeconomic environments \cite[p. 7]{filardo2008transparency}.

\paragraph{The European Central Bank Framework}
Within the Eurozone, the \gls{ecb} faces a communication challenge: it must address a fragmented, multi-national audience while maintaining a unified monetary stance. Since its establishment in 1999, the \gls{ecb}'s communication framework has progressively evolved to anchor expectations across diverse sovereign bond markets. This institutional evolution is most evident in its structured approach to policy announcements \cite{ECB_about}.

\subsection{Monetary Policy Statement}
\label{subsection:mps}
At the operational core of the \gls{ecb}'s communication strategy is the \gls{mps}, which is delivered during the press conference that directly follows the \gls{mp} meetings of the \gls{gc}. The council consists of the six members of the Executive Board (including the President and Vice-President) and the governors of the euro area national central banks, who meet to deliberate and set monetary policy parameters for the entire currency bloc. The \gls{ecb} currently schedules eight \gls{mp} meetings per year, operating on an approximate six-week cycle \cite{ITC_ECB_Schedule}.

\paragraph{Schedule of the Governing Council Meeting Day}
The organisation of the announcement day is standardised, reflecting the institution's strict commitment to predictability and the mitigation of market volatility. Starting from July 21, 2022, the \gls{ecb} modified its historical publication timeline to optimise the way financial markets process complex policy signals \cite{ECB2022_NewTimes}. The initial policy decision---usually a short press release detailing the exact interest rate parameters, deposit facility rates, and main structural changes to asset purchase programmes---is published at 14:15 CET. Shortly afterward, the press conference begins at 14:45 CET, during which the \gls{ecb} President and Vice-President present the detailed \gls{mps} and interact directly with the financial press. In parallel with the start of the \gls{qa} segment, the full text of the \gls{mps} is published on the \gls{ecb}'s digital platforms at 15:00 CET. At 17:00, European stock markets close independently of the \gls{ecb}, ending the day's reactions to its actions, as shown in Table~\ref{tab:meet_schedule}.

\begin{table}[ht]
    \centering
    \caption[Governing Council Meeting Day Schedule]{Governing Council Meeting Day Schedule.}
    \label{tab:meet_schedule}
    \begin{tabular}{@{}lr@{}}
        \toprule
        \textbf{Event}                         & \textbf{Scheduled time (CET)} \\ \midrule
        Meeting of the \gls{gc}                & 9:00 -- 12:00                 \\
        \gls{mp} announcement                  & 14:15                         \\
        Press conference (IS + \gls{qa}) & 14:45 -- 15:15                \\
        European stock market close            & 17:00                         \\ \bottomrule
    \end{tabular}
\end{table}


\paragraph{Structure of the Press Conference}
\label{sec:press_structure}
The press conference itself is divided into two distinct communication phases: the reading of the formally prepared \gls{is} and the subsequent unscripted \gls{qa} session. This structural division is not merely procedural; it has measurable consequences for financial markets, as the two sections serve different informational functions and generate distinct asset price dynamics.

The \gls{is} is a structured, carefully drafted document that represents the \gls{gc}'s consensus view \cite{ECB_WP2742}. It provides a linear overview of the economic assessment, inflation dynamics, and the evaluation of monetary and financial conditions. The language used in this document is deliberately precise, relying on specific syntactical formulations to signal the future policy path. This practice is commonly known as forward guidance: an explicit communication tool by which the central bank provides forward-looking indications about the future trajectory of its policy rates or asset purchases, thereby shaping market expectations beyond the current meeting \cite{blinder2008central}.

In contrast, the \gls{qa} session introduces a semi-structured interaction. Advanced textual analysis and structural topic modelling of historical transcripts demonstrate that the thematic content of the statement differs from the discourse during the \gls{qa} session. The topics dominating the statement are persistent, structural, and strongly focused on standardised macroeconomic aggregates. Conversely, the \gls{qa} often introduces new topics that were not addressed in the prepared text. Empirical research reveals that the first six most prevalent topics within a vast sample of press conference transcripts are predominantly driven by the \gls{qa} segments. This suggests that a significant share of the added value of the press conference lies in the \gls{qa} phase, which provides the public and market participants with detailed explanations regarding the \gls{ecb}'s decisions in real time \cite{ECB_WP2742}.

\subsection{Other Communication Channels}
While the \gls{mps} and the accompanying press conference represent the focal point of the \gls{ecb}'s communication strategy, they are complemented by a broader set of publications, reports, and public engagements. To continuously influence the yield curve, anchor long-term expectations, and fulfill its mandate of transparency, the \gls{ecb} applies several supplementary communication channels in the weeks between formal \gls{gc} meetings. This \emph{inter-meeting communication} is an effective mechanism for managing expectations between meetings and for addressing structural complexities and the mechanics of financial transmission. These details cannot be fully addressed in a single monthly press conference \cite{Istrefi2021ECB}.

\paragraph{The Economic Bulletin} is a comprehensive analytical publication that explains the reasons behind the \gls{mp} decisions of the \gls{ecb}'s Governing Council. Published eight times a year---generally two weeks after the corresponding \gls{mp} meeting---the Bulletin contextualises the Governing Council's decisions against the detailed backdrop of prevailing economic, financial, and monetary conditions \cite{press_to_speeches}.

\paragraph{The Monetary Policy Accounts} represent a significant transparency reform. The \gls{ecb} began publishing regular accounts of the Governing Council's \gls{mp} meetings in early 2015, aligning its institutional practices with those of the US \gls{fed} and the Bank of England. These accounts are typically released four weeks after the corresponding meeting \cite{MP_accounts}.

The documents offer a detailed overview of the Governing Council's internal deliberations. They outline the delicate balance of arguments, the range of risk assessments considered by various members, and the different perspectives on the economic outlook. By documenting the full range of analytical debates, the accounts help market participants understand the central bank's internal response function and the specific macroeconomic thresholds that might trigger future policy interventions \cite[p. 26--27]{CorbisieroLawton2021}.

As the \gls{ecb}'s communication channels have multiplied to foster transparency, the volume of textual data generated by the central bank has expanded substantially. While human analysts can interpret individual statements, the sheer scale and frequency of this qualitative information necessitate computational approaches to systematically extract the monetary policy signals. This demand for automated, large-scale text analysis motivates the application of \gls{nlp}.

\section{Natural Language Processing in Finance}
The growing availability of digital information has transformed the methodology of financial and macroeconomic research. Historically, empirical finance has relied almost exclusively on structured, quantitative data. However, language has always shaped financial markets, and unstructured textual data---ranging from earnings call transcripts to central bank \gls{mps}---contains critical information about the future. \gls{nlp} serves as a technological linkage, enabling researchers and practitioners to systematically extract structured, machine-readable signals from these vast corpora of unstructured text \cite{IBM_IE_2025}. In the context of \gls{mp}, \gls{nlp} algorithms are used to quantify the qualitative narratives of central bankers, capturing nuances in market psychology and institutional sentiment.

\subsection{Concept of Sentiment Analysis}
\label{sec:sentiment_concept}
Sentiment analysis is the computational extraction of subjective information and emotional states from text, with the central task of classifying a passage's underlying \textit{polarity} \cite{xing2020financial}. In its canonical formulation the polarity is mapped onto a positive--neutral--negative axis, where \textit{positive} denotes affirmative or optimistic framing and \textit{negative} denotes adverse or pessimistic framing.

Early applications in finance relied on dictionary-based approaches, which measure sentiment by aggregating the frequency of predefined terms. A common metric in this tradition is a normalised difference between positive and negative word counts. However, as noted by Loughran and McDonald \cite{LoughranMcDonald2011}, standard baseline dictionaries routinely misclassify common financial terminology; the word \textit{liability}, for example, is negative in general parlance but serves as a neutral, descriptive accounting term in financial filings. This mismatch motivated the development of domain-specific lexicons to accurately measure sentiment in corporate and macroeconomic contexts \cite{LoughranMcDonald2011}.

As deep learning matured, the methodology of sentiment quantification has shifted from counting isolated words to processing contextual probabilities. To facilitate econometric analysis, researchers transform these discrete class probabilities into a continuous \textit{Sentiment Score} ($S$). While the technology has evolved, the logic of aggregation remains similar to classic sentiment indices \cite{gentzkow2019text}. As defined in Equation~\ref{eq:sentiment}, the score is calculated as the difference between the probability assigned to the positive (or hawkish) class and its negative (or dovish) counterpart:

\begin{equation}
    \label{eq:sentiment}
    Score \;=\; P_{\text{positive}} \;-\; P_{\text{negative}},
\end{equation}

where $P \in [0, 1]$ represents the softmax probability assigned by the model to each respective label. This transformation maps the model's probabilistic output onto a linear scale bounded by $-1$ and $+1$. A value of zero signifies either a balanced or a neutrally classified passage. The two classifiers used in this thesis (introduced below) emit a binary positive/negative softmax: passages that the model assigns to a third \textit{neutral} class contribute zero by construction, attenuating the aggregate signal toward zero when a press conference consists predominantly of procedural or descriptive language.

The transition from static lexicons to transformer-based architectures, such as \textsc{bert} (Bidirectional Encoder Representations from Transformers), addresses the inability of bag-of-words models to capture context-dependent relationships. Unlike methods that process words in isolation, \textsc{bert} evaluates each token in relation to all others in a sentence, allowing for a nuanced treatment of syntax, negation, and complex economic narratives.

However, the linguistic requirements of central bank analysis nevertheless necessitate further specialisation. General-purpose models trained on Wikipedia or common literature struggle with the technical and coded lexicon used by monetary authorities \cite{BIS_WP1215_2024}. This limitation motivated the development of domain-adapted sentiment classifiers, of which two are employed in the present work. \textbf{FinBERT} \cite{araci2019finbert} is a BERT-base model further pre-trained on Reuters' TRC2 financial-news corpus and fine-tuned for sentiment on the Financial PhraseBank, and it represents the standard for general financial sentiment. \textbf{CentralBankRoBERTa} \cite{Pfeifer2023} is fine-tuned by Pfeifer and Marohl on a corpus of central bank speeches and reports and represents a further narrowing of the training distribution toward the institutional vocabulary of monetary authorities. Both models output a binary positive/negative polarity in the sense of Equation~\ref{eq:sentiment}; they differ in the domain of their pre-training corpus, not in their classification objective. Their differences are summarised in Table~\ref{tab:model_comparison}.

\begin{table}[ht]
    \centering
    \caption{Evolution of NLP models in financial and central bank contexts.}
    \label{tab:model_comparison}
    \begin{tabular}{@{}lll@{}}
        \toprule
        \textbf{Model} & \textbf{Training Corpus}                 & \textbf{Linguistic Nuance}  \\ \midrule
        BERT           & General (Wikipedia, BookCorpus)          & High (general context)      \\
        FinBERT        & Reuters TRC2 + Financial PhraseBank      & Domain-specific (financial) \\
        CBRoBERTa      & \gls{ecb}/\gls{fed} speeches and reports & Policy-specific (monetary)  \\ \bottomrule
    \end{tabular}
\end{table}

A conceptually distinct strand of the literature seeks to classify not the \emph{emotional} tone of central bank communications but their \emph{policy stance}: whether a passage signals a tightening (``hawkish''), an easing (``dovish'') or no change (``neutral'') in the monetary stance. Stance classification differs from sentiment analysis along two dimensions. First, its label space is operational rather than affective: a sentence such as \textit{``inflation is far too high''} carries a strongly negative emotional tone yet implies a hawkish policy stance, because elevated inflation is the standard precondition for tighter policy. Second, the construction of stance training sets requires manual annotation by domain experts using documented annotation guides, since the mapping from text to policy implication is not recoverable from sentiment cues alone. State-of-the-art stance classifiers for central bank discourse include the FOMC-trained RoBERTa-large model of Shah et al.\ \cite{shah2023trillion} and the multi-bank extension of Shah, Sukhani, Pardawala et al.\ \cite{shah2025wcb}, the latter of which provides a model fine-tuned specifically on \gls{ecb} communications.

These models are released under gated access on the Hugging Face Hub. As access could not be obtained within the timeframe of this study, the present analysis is restricted to sentiment classifiers, and the resulting signal should be interpreted as a measure of the \emph{affective tone} of \gls{ecb} discourse rather than its policy stance. Extending the pipeline to incorporate the stance dimension is identified in Chapter~\ref{ch:discussion} as a priority for future work.

\subsection{Zero-Shot Classification and the NLI Mechanism}
\label{sec:nlp_finance}
While the polarity score provides a measure of the tone of the discourse, it does not identify the subject of that tone. A hawkish signal regarding inflation has vastly different market implications than a hawkish remark regarding structural reforms. Sentiment analysis is therefore often complemented by thematic categorisation.

Traditionally, the thematic analysis of central bank communication was dominated by probabilistic generative models, most notably Latent Dirichlet Allocation (LDA) \cite{blei2003latent}. While useful for broad explorations, these ``bag-of-words'' approaches suffer from significant limitations, specifically an inability to capture semantic context and the requirement for users to pre-specify a fixed number of topics without a deeper understanding of linguistic nuance. To overcome these limitations, the field of financial NLP has increasingly shifted toward neural architectures using transformers and Natural Language Inference (NLI) frameworks \cite{yin2019benchmarking}.

This architectural shift enables \textit{zero-shot classification}: the technique allows a model to categorise text into user-defined labels without any task-specific supervised training. The procedure recasts classification as a logical relationship between a \textit{premise} (the original central bank passage) and a \textit{hypothesis} (a template representing the candidate label, such as \textit{``This passage discusses monetary policy.''}) \cite{yin2019benchmarking}.

To derive a confidence score for a specific category, the model's raw output logits for logical entailment and contradiction are processed through a softmax to yield the probability $P(c_i)$ for each candidate label $c_i$. In the implementation adopted in this work, the BART-large architecture pre-trained on the MultiNLI corpus \cite{lewis2019bart} provides the entailment backbone, and the neutral logit is discarded prior to normalisation:

\begin{equation}
    P(c_i) = \frac{\exp({l_{entail}})}{\exp({l_{entail}}) + \exp({l_{contra}})},
\end{equation}

where $l_{entail}$ and $l_{contra}$ represent the model's raw logits for entailment and contradiction, respectively. This mathematical framework allows the resulting thematic weight to reflect the relative strength of the logical link between the bank's statement and the defined hypothesis.

By combining domain-specific sentiment extraction with zero-shot thematic classification, the qualitative narratives of central bankers are transformed into structured, time-series data. However, to evaluate whether these linguistic signals actually influence financial markets, they must be benchmarked against observable economic variables. This requires an econometric framework capable of quantifying both the official policy stance and the high-frequency market reactions.

\section{Measuring Monetary Policy}
To fill the gap between qualitative NLP outputs and econometric modelling, it is essential to define the quantitative metrics that represent the monetary policy stance and market expectations. This section details the primary indicators used as targets for predictive modelling, ranging from official policy rates to structural market shocks. Table~\ref{tab:mp_indicators} summarises the four indicators that recur throughout the empirical chapters.

\paragraph{Official Policy Rates and the Shadow Rate}
The primary instrument of the \gls{ecb}'s monetary policy is the \gls{mro} rate, which sets the cost of liquidity for the banking system. While the \gls{mro} rate is the explicit signal of the policy stance, its effectiveness as a continuous variable is limited when interest rates reach the \gls{zlb}.

To measure the stance when nominal rates are constrained, researchers employ the \emph{Shadow Rate}. As proposed by \cite{krippner2013measuring} and operationalised for the euro area by Wu and Xia \cite{wu2016measuring}, the shadow rate is a model-based construct that translates unconventional policy measures, such as large-scale asset purchases and forward guidance, into a single interest rate equivalent. Unlike the \gls{mro}, the shadow rate can move into negative territory, providing a more granular metric of monetary conditions during periods of intensive central bank intervention.

\paragraph{Market Expectations and OIS Volatility}
Real-time market reactions to central bank communication are primarily captured using \gls{ois}. An \gls{ois} is a financial derivative where a fixed interest rate is exchanged for a floating overnight rate (specifically the  €STR in the Eurozone) over the duration of the contract. Since there is no exchange of principal and the floating leg is based on an overnight rate, \gls{ois} rates are considered a ``risk-free'' proxy for the market's expectation of the future path of monetary policy.

In the context of the \emph{EA-MPD} dataset \cite{altavilla2019measuring}, the focus is on the \emph{change} in these rates ($\Delta \text{OIS}$) during the press conference window. The dataset provides a term structure of expectations ranging from one month to ten years. Each segment serves a distinct analytical purpose: the short end (1M to 6M) reflects expectations regarding near-term policy rates, the mid-curve (1Y to 2Y) captures the effectiveness of forward guidance, and the long end (5Y to 10Y) acts as a gauge for long-term inflation anchoring.

Beyond the directional shifts in rates, their intraday volatility during the announcement serves as a critical proxy for market uncertainty. A divergence between the \gls{is} and the spontaneous \gls{qa} session---often referred to as a ``communication mismatch''---typically manifests as heightened volatility. This indicates that the market is struggling to price in conflicting signals, making uncertainty a quantifiable dimension of the central bank's communication effectiveness.

\begin{table}[ht]
    \centering
    \caption{Key Monetary Policy Indicators used as empirical targets.}
    \label{tab:mp_indicators}
    \begin{tabular}{@{}ll@{}}
        \toprule
        \textbf{Indicator}        & \textbf{Economic Significance}                                        \\ \midrule
        \gls{mro} Rate            & Direct policy signal and cost of liquidity.                           \\
        Shadow Rate               & Policy stance proxy in \gls{zlb} environments \cite{wu2016measuring}. \\
        \gls{ois} Rates (1M--10Y) & Forward-looking expectations across the yield curve.                  \\
        $PC1$ (Shocks)            & Aggregated policy surprise \cite{jarocinski2020deconstructing}.       \\ \bottomrule
    \end{tabular}
\end{table}

\paragraph{Structural Shocks and Principal Component Analysis}
The ``pure'' impact of a policy announcement is often clouded by the information effect, where markets react to the central bank's economic outlook rather than the policy change itself. To isolate the intended policy signal, \cite{jarocinski2020deconstructing} and \cite{altavilla2019measuring} utilise structural decomposition. By observing the co-movement of interest rates and stock prices, researchers separate these shocks into ``pure monetary'' and ``information'' components. To simplify the multi-dimensional nature of interest rate moves across the yield curve, Principal Component Analysis (PCA) is applied to the high-frequency $\Delta \text{OIS}$ vector. The First Principal Component ($PC1$) is extracted to represent the overarching shift in expectations, as shown in Equation~\ref{eq:pc1}:

\begin{equation}
    \label{eq:pc1}
    PC1_t = \sum_{j=1}^{n} \omega_j \cdot \text{Shock}_{j,t},
\end{equation}

where $\omega_j$ represents the PCA loading for the structural shock at maturity $j$. This $PC1$ serves as a robust target variable for the machine-learning models of Chapter~\ref{ch:results}, as it captures the dominant policy surprise while filtering out idiosyncratic noise from individual maturities.

\paragraph{Chapter Synthesis}
The institutional structure of central bank communication, the computational extraction of sentiment via \gls{nlp}, and the quantitative measurement of monetary policy shocks form the basis of this thesis. By mapping the qualitative shifts in the \gls{ecb}'s discourse onto the structural changes in the \gls{ois} yield curve, the methodology of Chapter~\ref{ch:data} details the empirical framework used to isolate the trading value of monetary policy signals, and the empirical findings of Chapter~\ref{ch:results} evaluate it against the financial-market benchmarks defined above.